\documentclass[11pt]{article}
\usepackage{amsmath,amsthm,amssymb,amsfonts,geometry,hyperref}
\usepackage[ruled,vlined]{algorithm2e}
\geometry{margin=1in}

\newtheorem{definition}{Definition}[section]
\newtheorem{remark}[definition]{Remark}
\newtheorem{lemma}[definition]{Lemma}
\newtheorem{theorem}[definition]{Theorem}
\newtheorem{proposition}[definition]{Proposition}
\newtheorem{corollary}[definition]{Corollary}
\hypersetup{colorlinks=true,linkcolor=blue,citecolor=blue,urlcolor=blue}

\title{\textbf{Validated Spectral Analysis of the Gauss Transfer Operator\\
and the Gauss--Kuzmin--Wirsing Law}}
\author{Isaia Nisoli}
\date{}

\begin{document}
\maketitle

\begin{abstract}
We develop a rigorous Hardy-space framework for the Gauss transfer operator
\[
(L f)(w)=\sum_{n=1}^{\infty}\frac{1}{(w+n)^2}\,f\!\left(\frac{1}{w+n}\right),
\]
acting from \(H^2(D_1)\) to \(H^2(D_{3/2})\).
We prove compactness, derive explicit norm and truncation estimates,
and implement a validated numerical scheme to enclose multiple eigenvalues
and eigenvectors of \(L\).
This allows a certified computation of the Gauss--Kuzmin--Wirsing constant
and a complete spectral expansion for continued-fraction digit distributions,
answering Gauss's 1812 question with rigorous error bounds.
\end{abstract}

\section{Introduction and Main Results}

The Gauss map \(T(x)=\{1/x\}\) on \((0,1]\) generates continued fractions.
Gauss asked for the probability
\[
F_n(x)=\mathbb P(a_n\le x)
\]
that the \(n\)-th digit \(a_n\) of the continued-fraction expansion is at most \(x\),
if \(x_0\) is uniformly distributed in \([0,1]\).
This distribution is governed by iterates of the transfer operator
\[
(L f)(x)=\sum_{n\ge1}\frac{1}{(x+n)^2}f\!\left(\frac{1}{x+n}\right).
\]

We work in Hardy spaces \(H^2(D_r)\), with \(D_r=\{|w-1|<r\}\),
derive explicit analytic bounds, and rigorously enclose the spectrum of \(L\).

\paragraph{Main results.}
\begin{itemize}
\item \(L:H^2(D_1)\to H^2(D_{3/2})\) is compact with
approximation numbers \(a_m(L)\le C(2/3)^m\).
\item Certified enclosures of eigenvalues and eigenvectors,
including rigorous bounds on the Gauss--Kuzmin--Wirsing constant.
\item Spectral expansion
\[
F_n(x)=\frac{\ln(1+x)}{\ln2}
+\sum_{j\ge2}c_j\lambda_j^{\,n}H_j(x),
\quad H_j(x)=\!\int_0^x\phi_j(t)\,dt,
\]
with error \(O(\rho^n)\).
\end{itemize}

\section{Analytic Framework}\label{sec:analytic-bounds}

\begin{definition}[Hardy spaces]
For \(r>0\), \(H^2(D_r)\) denotes the space of holomorphic functions on
\(D_r=\{|w-1|<r\}\) with norm
\[
\|f\|_{H^2(D_r)}^2=\sum_{k\ge0}|c_k|^2 r^{2k},
\qquad f(w)=\sum_{k\ge0}c_k(w-1)^k.
\]
\end{definition}

\begin{lemma}[Geometry of the branch maps]
Let \(\tau_n(w)=1/(w+n)\).
Then for all \(n\ge1\), \(\tau_n(D_{3/2})\subset D_1\).
\end{lemma}

\begin{proof}
For \(w\in D_{3/2}\), \(\Re w>-\tfrac12\).
Compute
\[
|\tau_n(w)-1|=\frac{|w+n-1|}{|w+n|},
\qquad
|w+n|^2-|w+n-1|^2=2\Re w+(2n-1).
\]
Since \(\Re w>-\tfrac12\), we get \(2\Re w+(2n-1)>2n-2\ge0\),
with equality only for \(n=1\) and \(w=-\tfrac12\) on \(\partial D_{3/2}\).
Hence \(|w+n-1|<|w+n|\) and \(|\tau_n(w)-1|<1\), proving the claim.
\end{proof}

\begin{proposition}
\label{prop:branch-H2}
For each $n\in\mathbb N$, the composition operator
\[
C_{\tau_n}:H^2(D_1)\to H^2(D_{3/2}),\qquad C_{\tau_n}f=f\circ\tau_n,
\]
satisfies
\[
\|C_{\tau_n}\|_{H^2(D_1)\to H^2(D_{3/2})}
\ \le\
\Big(\frac{1+|\phi_n(0)|}{1-|\phi_n(0)|}\Big)^{1/2}
=\sqrt{2n+1},
\]
where $\phi_n=T_1^{-1}\circ\tau_n\circ T_{3/2}$ with $T_r(\zeta)=1+r\zeta$.
\end{proposition}

\begin{proof}
Let $U_r:H^2(D_r)\to H^2(\mathbb D)$ be the unitary operator $(U_r f)(\zeta)=f(T_r(\zeta))$.
Then
\[
U_{3/2}\, C_{\tau_n}\, U_1^{-1} \;=\; C_{\phi_n},
\qquad
\phi_n=T_1^{-1}\circ\tau_n\circ T_{3/2},
\]
so $\|C_{\tau_n}\|=\|C_{\phi_n}\|$. A direct computation gives
\[
\phi_n(\zeta)=\frac{-\,n-\tfrac32\zeta}{\,n+1+\tfrac32\zeta\,},
\qquad
\phi_n(0)=-\frac{n}{n+1}.
\]
By \cite[Thm.~3.6]{CowenMacCluer2019}, for $H^2(\mathbb D)$ we have
\[
\|C_{\phi}\|_{H^2(\mathbb D)}\le\big(\tfrac{1+|\phi(0)|}{1-|\phi(0)|}\big)^{1/2}.
\]
Applying this to $\phi=\phi_n$ yields
\[
\|C_{\tau_n}\|_{H^2(D_1)\to H^2(D_{3/2})}
=\|C_{\phi_n}\|_{H^2(\mathbb D)}
\le \Big(\frac{1+\frac{n}{n+1}}{1-\frac{n}{n+1}}\Big)^{1/2}
=\sqrt{2n+1}.
\]
\end{proof}



\begin{lemma}\label{lem:weight-bound}
For \(w\in D_{3/2}\),
\[
|(w+n)^{-2}|\le(n-\tfrac12)^{-2}.
\]
\end{lemma}

\begin{proof}
The minimum of \(|w+n|\) on \(\overline{D_{3/2}}\) is at \(w=-\tfrac12\):
\(|w+n|\ge n-\tfrac12\).
\end{proof}

\begin{theorem}
\label{thm:L1-H2-bound}
We have
\[
\|L\|_{H^2(D_1)\to H^2(D_{3/2})}
\ \le\
\sum_{n\ge1}\frac{\sqrt{2n+1}}{(n-\tfrac12)^2}
\ \le\
\sqrt{6}\,2^{3/2}\!\bigl(1-2^{-3/2}\bigr)\,\zeta\!\Big(\tfrac32\Big)
=:C_{2}.
\]
\end{theorem}
\begin{proof}
From Lemma \ref{lem:weight-bound} and Proposition~\ref{prop:branch-H2},
\[
\|L f\|_{H^2(D_1)\to H^2(D_{3/2})}
\le \sum_{n\ge1}\frac{1}{(n-\tfrac12)^2}\,\|C_{\tau_n}\|\,\|f\|
\le \sum_{n\ge1}\frac{\sqrt{2n+1}}{(n-\tfrac12)^2}\,\|f\|.
\]
Use the elementary comparison
\(
\frac{2n+1}{n-\tfrac12}\le 6
\)
to get \(\sqrt{2n+1}\le \sqrt6\,(n-\tfrac12)^{1/2}\).

Now, we observe that 
\[
\sum_{n\geq 1}(n-\tfrac12)^{-s} = \sum_{m\geq 0}(m+\tfrac12)^{-s} = 2^s\sum_{m\geq 0}(2m+1)^{-s},
\]
and that 
\[
\sum_{m=0}^\infty (2m+1)^{-s}
= \zeta(s)-\sum_{m=1}^\infty (2m)^{-s}
= \zeta(s)-2^{-s}\zeta(s)
= (1-2^{-s})\,\zeta(s), \qquad \Re s>1.
\]
Thus:
\[
\sum_{n\geq 1}(n-\tfrac12)^{-s} = 2^s (1-2^{-s})\,\zeta(s) = (2^s - 1)\,\zeta(s).
\]
Hence
\[
\sum_{n\ge1}\frac{\sqrt{2n+1}}{(n-\tfrac12)^2}
\le \sqrt6\sum_{n\ge1}(n-\tfrac12)^{-3/2}
=\sqrt6\,\!\bigl(2^{3/2}-1\bigr)\,\zeta(\tfrac32).
\]
\end{proof}
\begin{remark}
% A rigorous numerical evaluation gives an upper bound of 
% \[
% C_2 \leq 11.7000807085501779306421385006549
% \]
We observe that
\[
\frac{\sqrt{2n+1}}{(n-\tfrac12)^2}
=\sqrt{\frac{2n+1}{\,n-\tfrac12\,}}\,(n-\tfrac12)^{-3/2},
\qquad
\sup_{n\ge N}\frac{2n+1}{n-\tfrac12}
=2+\frac{2}{N-\tfrac12}.
\]
Hence, for any $N\ge1$,
\[
\|L\|_{H^2(D_1)\to H^2(D_{3/2})}
\le \sum_{n=1}^{N-1}\frac{\sqrt{2n+1}}{(n-\tfrac12)^2}
\;+\;\sqrt{\,2+\frac{2}{N-\tfrac12}\,}\,
\sum_{n\ge N}(n-\tfrac12)^{-3/2}.
\]
Using the Hurwitz zeta,
\(\displaystyle \sum_{n\ge N}(n-\tfrac12)^{-3/2}=\zeta\!\left(\tfrac32,\,N-\tfrac12\right)\).

In table \ref{tab:C2-bounds} are rigorously computed upper bounds for $C_2$ as a function of the splitting $N$, $N=1$ is the base case from Theorem~\ref{thm:L1-H2-bound},
using Arb, with maximum enclosure radius bounded by $5\times10^{-34}$.

\end{remark}

\begin{table}[h]
\centering
\begin{tabular}{r r}
\hline
$N$ & $C_2(N)$ \\ \hline
1 & 11.7000807086 \\
2 & 10.4849507042 \\
3 & 10.2709840576 \\
4 & 10.1907129342 \\
5 & 10.1506945125 \\
10 & 10.0893118365 \\
100 & 10.0590444185 \\
1000 & 10.0581271764 \\
10000 & 10.0580982915 \\
\hline
\end{tabular}
\caption{Refined upper bounds for \(C_2\) as a function of the splitting \(N\).}\label{tab:C2-bounds}
\end{table}


\section{Truncation bounds and approximation numbers}


Compare the analogous estimate in the $H^\infty$–ellipse setting: \cite[Thm.~3.3(ii)]{BandtlowSlipantschuk2020}. 
See also the published variant using the same factorization idea in \cite{BandtlowPohlSchickWeisse2021}.

\begin{lemma}
\label{lem:H2-tail}
Let \(0<r<R\) and let \(f(w)=\sum_{n\ge0} a_n (w-1)^n\) be holomorphic on \(D_R=\{|w-1|<R\}\) with
\[
\|f\|_{H^2(D_R)}^2=\sum_{n\ge0} |a_n|^2 R^{2n}<\infty.
\]
Let \(\Pi_k f:=\sum_{n=0}^k a_n (w-1)^n\) be the degree-\(k\) truncation, and let
\(J:H^2(D_R)\to H^2(D_r)\) be the restriction map (inclusion at the level of Banach spaces). Then
\[
\|\,Jf-\Pi_k f\,\|_{H^2(D_r)}
\;\le\;
\Bigl(\frac{r}{R}\Bigr)^{k+1}\,\|f\|_{H^2(D_R)}.
\]
In particular, for \(r=1\) and \(R=\tfrac32\),
\[
\|\,Jf-\Pi_k f\,\|_{H^2(D_1)}
\;\le\;
\Bigl(\frac{2}{3}\Bigr)^{k+1}\,\|f\|_{H^2(D_{3/2})}.
\]
\end{lemma}

\begin{proof}
First, \(\|J\|_{H^2(D_R)\to H^2(D_r)}=1\): indeed,
\(\sum |a_n|^2 r^{2n}\le \sum |a_n|^2 R^{2n}\) and equality holds for constants.
By the definition of the \(H^2(D_r)\) norm,
\[
\|Jf-\Pi_k f\|_{H^2(D_r)}^2
=\sum_{n>k} |a_n|^2 r^{2n}
=\sum_{n>k} |a_n|^2 R^{2n}\,\Bigl(\frac{r}{R}\Bigr)^{2n}.
\]
Since \(0<\frac{r}{R}<1\) and \(n>k\), we have
\(\bigl(\frac{r}{R}\bigr)^{2n}\le \bigl(\frac{r}{R}\bigr)^{2(k+1)}\). Therefore
\[
\|Jf-\Pi_k f\|_{H^2(D_r)}^2
\le \Bigl(\frac{r}{R}\Bigr)^{2(k+1)} \sum_{n>k} |a_n|^2 R^{2n}
\le \Bigl(\frac{r}{R}\Bigr)^{2(k+1)} \sum_{n\ge0} |a_n|^2 R^{2n}
= \Bigl(\frac{r}{R}\Bigr)^{2(k+1)} \|f\|_{H^2(D_R)}^2,
\]
and taking square roots gives the claim.
\end{proof}

\begin{theorem}
\label{thm:trunc-H2-correct}
Let $T:=R\,L\,S: H^2(D_{3/2})\to H^2(D_1)$, where $S$ and $R$ are the restriction maps
$H^2(D_{3/2})\to H^2(D_1)$. Let $\Pi_k$ be the degree-$k$ truncation on $H^2(D_{3/2})$
and $\pi_K$ the degree-$K$ truncation on $H^2(D_{3/2})$ (acting on the \emph{range} of $L$).
Then, with $C_2:=\|L\|_{H^2(D_1)\to H^2(D_{3/2})}$ (e.g. $C_2\le\sqrt6\,2^{3/2}(1-2^{-3/2})\zeta(\tfrac32)$),
\[
\big\|\,T - R\,\pi_K\,L\,\Pi_k\,S\,\big\|_{H^2(D_{3/2})\to H^2(D_1)}
\ \le\
C_2\!\left[\left(\frac{2}{3}\right)^{K+1}
+\left(\frac{2}{3}\right)^{k+1}\right].
\]
\end{theorem}

\begin{proof}
Decompose
\[
T - R\,\pi_K\,L\,\Pi_k\,S
= R\,(I-\pi_K)\,L\,S \;+\; R\,\pi_K\,L\,(I-\Pi_k)\,S
=: A_K + B_k .
\]
Since $\|R\|=\|S\|=1$ and $\|L\|=C_2$, it suffices to bound the truncation tails after
restricting from radius $3/2$ to radius $1$.

For $A_K$ (target-side tail):
by Lemma~\ref{lem:H2-tail} with $(r,R)=(1,3/2)$ applied to $g\in H^2(D_{3/2})$,
\[
\|R\,(I-\pi_K)g\|_{H^2(D_1)}
\ \le\ \Big(\frac{2}{3}\Big)^{K+1}\,\|g\|_{H^2(D_{3/2})}.
\]
Hence
\[
\|A_K\|
=\sup_{\|f\|_{H^2(D_{3/2})}=1}
\|R\,(I-\pi_K)\,L\,S f\|_{H^2(D_1)}
\ \le\ \Big(\frac{2}{3}\Big)^{K+1}\,\|L\|
= \Big(\frac{2}{3}\Big)^{K+1} C_2.
\]

For $B_k$ (domain-side tail):
apply Lemma~\ref{lem:H2-tail} to $f\in H^2(D_{3/2})$ to get
\[
\|S f - S\Pi_k f\|_{H^2(D_1)}
\ \le\ \Big(\frac{2}{3}\Big)^{k+1}\,\|f\|_{H^2(D_{3/2})}.
\]
Therefore
\[
\|B_k\|
\le \|R\|\,\|\pi_K\|\,\|L\|\,\|S(I-\Pi_k)\|
\le C_2\,\Big(\frac{2}{3}\Big)^{k+1}.
\]

Combining the two bounds yields the claim.
\end{proof}


% \begin{corollary}
% \label{cor:approxnums}
% Let $T:=R\,L\,S: H^2(D_{3/2})\to H^2(D_1)$, where $S$ and $R$ are the
% restriction maps $H^2(D_{3/2})\to H^2(D_1)$ (both of norm $1$).
% Let $a_m(T)$ denote the $m$-th approximation number of $T$, i.e.
% \[
% a_m(T)\ :=\ \inf\{\ \|T-F\|\ :\ \mathrm{rank}(F)<m\ \}.
% \]
% Then, with $C_2:=\|L\|_{H^2(D_1)\to H^2(D_{3/2})}$, one has the
% geometric bound
% \[
% \boxed{\qquad a_m(T)\ \le\ 2\,C_2\ \Bigl(\frac{2}{3}\Bigr)^{m}
% \qquad (m\ge1). \qquad}
% \]
% In particular, $T$ is compact; moreover $\sum_{m\ge1} a_m(T) <\infty$, so
% $T$ is trace class on $H^2(D_1)$.
% \end{corollary}

% \begin{proof}
% For integers $K,k\ge0$ define the finite-rank operator
% \[
% T_{K,k}\ :=\ R\,\pi_K\,L\,\Pi_k\,S.
% \]
% Its range is contained in the $(K+1)$-dimensional subspace
% $\mathrm{span}\{(w-1)^n\}_{n=0}^K$ of $H^2(D_1)$, hence
% $\mathrm{rank}(T_{K,k})\le K+1$. Also, $\Pi_k$ has range of dimension
% $k+1$, so $\mathrm{rank}(T_{K,k})\le k+1$. Thus
% \[
% \mathrm{rank}(T_{K,k})\ \le\ \min\{K+1,k+1\}.
% \]

% By the truncation inequality (Theorem~\ref{thm:trunc-H2-correct}),
% \[
% \|T-T_{K,k}\|
% \ \le\ C_2\Bigl[\Bigl(\tfrac{2}{3}\Bigr)^{K+1}
%                  +\Bigl(\tfrac{2}{3}\Bigr)^{k+1}\Bigr].
% \]
% Given $m\ge1$, choose $K=k=m-1$. Then $\mathrm{rank}(T_{K,k})\le m$ and
% \[
% \|T-T_{K,k}\|\ \le\ 2\,C_2\Bigl(\tfrac{2}{3}\Bigr)^{m}.
% \]
% Taking the infimum over all rank-$<m$ operators yields
% $a_m(T)\le 2\,C_2(2/3)^{m}$.

% Compactness follows since $a_m(T)\to0$. Furthermore,
% $\sum_{m\ge1}a_m(T)\le 2C_2\sum_{m\ge1}(2/3)^{m}<\infty$, so $T$ is
% trace class.
% \end{proof}

\begin{corollary}[One–sided finite–rank approximations of $L_{1}$ and $L_{3/2}$]
\label{cor:one-sided-approxs}
Let $L:H^2(D_1)\to H^2(D_{3/2})$ be bounded with $\|L\|=C_2$, and set
$L_{1}=S\,L:H^2(D_1)\to H^2(D_1)$ and $L_{3/2}=L\,R:H^2(D_{3/2})\to H^2(D_{3/2})$.

\smallskip
\noindent\textup{(i) Range truncation for $L_{1}$.}
For $K\ge0$ define the finite–rank operator
\[
(L_{1})_{K}:=\Pi_{K}\,L_{1}=\Pi_{K}\,S\,L\quad(H^2(D_1)\to H^2(D_1)).
\]
Then
\[
\|\,L_{1}-(L_{1})_{K}\,\|_{H^2(D_1)\to H^2(D_1)}
\ \le\ C_2\Bigl(\frac{2}{3}\Bigr)^{K+1}.
\]

\noindent\textup{(ii) Domain truncation for $L_{3/2}$.}
For $k\ge0$ define the finite–rank operator
\[
(L_{3/2})_{k}:=L\,\Pi_{k}\,R\quad(H^2(D_{3/2})\to H^2(D_{3/2})).
\]
Then
\[
\|\,L_{3/2}-(L_{3/2})_{k}\,\|_{H^2(D_{3/2})\to H^2(D_{3/2})}
\ \le\ C_2\Bigl(\frac{2}{3}\Bigr)^{k+1}.
\]
\end{corollary}

\begin{proof}
(i) For $f\in H^2(D_1)$, $L_{1}f=S(Lf)$ is the restriction to $D_1$ of a function in $H^2(D_{3/2})$
with norm $\le C_2\|f\|_{H^2(D_1)}$. Apply Lemma~\ref{lem:H2-tail} with $(r,R)=(1,3/2)$ to
$g:=L f\in H^2(D_{3/2})$:
\[
\|\,L_{1}f-\Pi_{K}L_{1}f\,\|_{H^2(D_1)}=\|\,Sg-\Pi_K Sg\,\|_{H^2(D_1)}
\le \Bigl(\tfrac{2}{3}\Bigr)^{K+1}\|g\|_{H^2(D_{3/2})}
\le C_2\Bigl(\tfrac{2}{3}\Bigr)^{K+1}\|f\|_{H^2(D_1)}.
\]
Taking the supremum gives the claim.

(ii) For $F\in H^2(D_{3/2})$,
\[
\|\,L_{3/2}F-(L_{3/2})_{k}F\,\|_{H^2(D_{3/2})}
=\|\,L(RF-\Pi_k RF)\,\|_{H^2(D_{3/2})}
\le C_2\,\|\,RF-\Pi_k RF\,\|_{H^2(D_1)}.
\]
Apply Lemma~\ref{lem:H2-tail} with $(r,R)=(1,3/2)$ to $F$ to obtain
$\|RF-\Pi_k RF\|_{H^2(D_1)}\le (2/3)^{k+1}\|F\|_{H^2(D_{3/2})}$, which yields the bound.
\end{proof}

\begin{proposition}
\label{prop:L-on-each-space}
Let $L:H^2(D_1)\to H^2(D_{3/2})$ be bounded with $\|L\|=C_2$, and set
$L_{1}=S\,L:H^2(D_1)\to H^2(D_1)$ and $L_{3/2}=L\,R:H^2(D_{3/2})\to H^2(D_{3/2})$.
Then:
\begin{enumerate}
\item $L_{1}$ and $L_{3/2}$ are compact.
\item For all $m\ge1$,
\[
a_m(L_{1})\ \le\ C_2\Bigl(\tfrac{2}{3}\Bigr)^{m-1},
\qquad
a_m(L_{3/2})\ \le\ C_2\Bigl(\tfrac{2}{3}\Bigr)^{m-1}.
\]
\end{enumerate}
\end{proposition}
\begin{proof}
By Corollary~\ref{cor:one-sided-approxs}, for every $K\ge0$ the rank-$\le K+1$
operator $(L_{1})_{K}=\Pi_K S L$ satisfies
$\|L_{1}-(L_{1})_{K}\|\le C_2(2/3)^{K+1}$; hence $\|L_{1}-(L_{1})_{K}\|\to0$,
so $L_{1}$ is compact. The same argument with $(L_{3/2})_{k}=L\Pi_k R$
gives compactness of $L_{3/2}$.

For the approximation numbers, fix $m\ge1$. If $m=1$, then
$a_1(L_{1})\le\|L_{1}\|\le C_2$ and similarly for $L_{3/2}$, which matches the
stated bound. If $m\ge2$, choose $K=k=m-2$ so that
$\operatorname{rank}((L_{1})_{K})\le m-1$ and $\operatorname{rank}((L_{3/2})_{k})\le m-1$.
Then
\[
a_m(L_{1})\ \le\ \|L_{1}-(L_{1})_{m-2}\|
\ \le\ C_2\Bigl(\tfrac{2}{3}\Bigr)^{(m-2)+1}
\ =\ C_2\Bigl(\tfrac{2}{3}\Bigr)^{m-1},
\]
and the same estimate holds for $a_m(L_{3/2})$.
\end{proof}

\begin{remark}
An alternative proof of the compactness of $J_{R\to r}$ is to use the normal–family
(Montel) theorem: the unit ball of $H^2(D_R)$ is locally bounded, hence normal; any
bounded sequence has a subsequence converging uniformly on $\overline{D_r}$, and the
coefficient tail estimate of Lemma~\ref{lem:H2-tail} upgrades this to convergence in
the $H^2(D_r)$–norm, yielding compactness.
\end{remark}

\begin{corollary}[Eigenvalue decay and trace class for $L_{1}$ and $L_{3/2}$]
\label{cor:eig-decay-traceclass}
Let $L_{1}=SL:H^2(D_1)\to H^2(D_1)$ and $L_{3/2}=LR:H^2(D_{3/2})\to H^2(D_{3/2})$
be as in Proposition~\ref{prop:L-on-each-space}. Enumerate the eigenvalues
$\{\lambda_m(\cdot)\}_{m\ge1}$ in nonincreasing order of modulus, with multiplicity.
Then, with $C_2=\|L\|_{H^2(D_1)\to H^2(D_{3/2})}$,
\[
|\lambda_m(L_{1})|\ \le\ C_2\Bigl(\tfrac{2}{3}\Bigr)^{m-1},
\qquad
|\lambda_m(L_{3/2})|\ \le\ C_2\Bigl(\tfrac{2}{3}\Bigr)^{m-1}
\qquad(m\ge1).
\]
In particular, both $L_{1}$ and $L_{3/2}$ are trace class and satisfy the trace–norm bounds
\[
\|L_{1}\|_{S_1}\ \le\ 3\,C_2,
\qquad
\|L_{3/2}\|_{S_1}\ \le\ 3\,C_2,
\]
hence $\displaystyle \sum_{m\ge1}|\lambda_m(L_{1})|<\infty$ and
$\displaystyle \sum_{m\ge1}|\lambda_m(L_{3/2})|<\infty$.
\end{corollary}

\begin{proof}
By Proposition~\ref{prop:L-on-each-space}, the approximation numbers satisfy
$a_m(L_{1}),\,a_m(L_{3/2})\le C_2(2/3)^{m-1}$ for $m\ge1$. On a Hilbert space,
$a_m(T)=s_m(T)$ equals the $m$th singular value, and Weyl’s inequality gives
$|\lambda_m(T)|\le s_m(T)$ (see Simon~\cite[Thm.~1.33]{Simon05}). Thus
$|\lambda_m(L_{1})|,|\lambda_m(L_{3/2})|\le C_2(2/3)^{m-1}$.

Moreover, $\sum_{m\ge1}s_m(L_{1})\le C_2\sum_{m\ge1}(2/3)^{m-1}=3C_2$, and similarly for
$L_{3/2}$. Hence $L_{1}$ and $L_{3/2}$ are trace class with the stated trace–norm bounds
(see Simon~\cite[Thm.~1.35]{Simon05}). Finally, since $|\lambda_m|\le s_m$, the series of
absolute eigenvalues is summable.
\end{proof}

\begin{remark}
\label{rem:Ak-norms}
Let $A_k$ denote the Galerkin compression of either endomorphism
$L_1:=SL:H^2(D_1)\to H^2(D_1)$ or $L_{3/2}:=LR:H^2(D_{3/2})\to H^2(D_{3/2})$
onto $\mathcal V_N=\mathrm{span}\{(w-1)^n\}_{n=0}^{N-1}$, with orthogonal
projection taken in the $H^2(D_r)$ inner product. Then
\[
\|A_k\|_{(r)}\ \le\ \|L_r\|\ \le\ C_2,\qquad r\in\{1,\tfrac32\}.
\]
Indeed, $\|L_1\|\le\|S\|\,\|L\|=C_2$ and $\|L_{3/2}\|\le\|L\|\,\|R\|=C_2$
(since $\|R\|=\|S\|=1$), and compressions do not increase the operator norm.
\end{remark}

\section{Spectral stability and projections from the discretization}
\label{sec:spectral-stability}

Recall the endomorphisms
\[
L_{1}:=S\,L:H^2(D_1)\to H^2(D_1),
\qquad
L_{3/2}:=L\,R:H^2(D_{3/2})\to H^2(D_{3/2}),
\]
and their one–sided finite–rank approximants from Corollary~\ref{cor:one-sided-approxs}:
\[
(L_{1})_{K}=\Pi_{K} S L,\qquad
(L_{3/2})_{k}=L \Pi_{k} R,
\]
with
\[
\|L_{1}-(L_{1})_{K}\| \le \varepsilon_{K}:=C_2\Bigl(\tfrac{2}{3}\Bigr)^{K+1},
\qquad
\|L_{3/2}-(L_{3/2})_{k}\| \le \delta_{k}:=C_2\Bigl(\tfrac{2}{3}\Bigr)^{k+1}.
\]

We now record a quantitative ``resolvent bridge'' and its consequences for enclosing
eigenvalues and approximating spectral projections.

\begin{proposition}\label{prop:ev-norm-bound}
Let $(\lambda,\varphi)$ be a (nonzero) eigenpair of $L$ with
$\varphi\in H^2(D_1)\cap H^2(D_{3/2})$ and $L\varphi=\lambda\varphi$.
Then
\[
\boxed{\qquad
\|\varphi\|_{H^2(D_{3/2})}
\ \le\ \frac{C_2}{|\lambda|}\,\|\varphi\|_{H^2(D_1)}.
\qquad}
\]
In particular, under the normalization $\|\varphi\|_{H^2(D_1)}=1$,
\[
\|\varphi\|_{H^2(D_{3/2})}\ \le\ \frac{C_2}{|\lambda|}.
\]
\end{proposition}

\begin{proof}
Since $L\varphi=\lambda\varphi$ and $L:H^2(D_1)\to H^2(D_{3/2})$,
\[
\|\varphi\|_{H^2(D_{3/2})}
=\frac{1}{|\lambda|}\,\|L\varphi\|_{H^2(D_{3/2})}
\le \frac{C_2}{|\lambda|}\,\|\varphi\|_{H^2(D_1)}.
\]
\end{proof}


\begin{lemma}[Resolvent bridge]\label{lem:resolvent-bridge}
Let $A,B$ be bounded operators on a Banach space and fix $z\in\mathbb{C}$ with
$z\in\rho(B)$ (the resolvent set). If
\[
\alpha(z):=\|(A-B)R_B(z)\|\;<\;1,
\]
then $z\in\rho(A)$ and
\begin{equation}\label{eq:resolvent-bridge-bounds}
\|R_A(z)\|\ \le\ \frac{\|R_B(z)\|}{1-\alpha(z)},
\qquad
\|R_A(z)-R_B(z)\|
\ \le\ \frac{\|R_B(z)\|^2\,\|A-B\|}{1-\alpha(z)}.
\end{equation}
Moreover, if $\Gamma$ is a positively oriented Jordan curve with $\Gamma\subset\rho(B)$
and $\sup_{z\in\Gamma}\alpha(z)<1$, then $\Gamma\subset\rho(A)$ and the Riesz projections
\[
P_A(\Gamma)=\frac{1}{2\pi i}\int_\Gamma R_A(z)\,dz,\qquad
P_B(\Gamma)=\frac{1}{2\pi i}\int_\Gamma R_B(z)\,dz
\]
are both well defined and satisfy
\begin{equation}\label{eq:proj-diff}
\|P_A(\Gamma)-P_B(\Gamma)\|
\ \le\ \frac{|\Gamma|}{2\pi}\ \sup_{z\in\Gamma}
\frac{\|R_B(z)\|^2\,\|A-B\|}{1-\alpha(z)},
\end{equation}
and $A$ and $B$ have the same total algebraic multiplicity of eigenvalues inside $\Gamma$.
\end{lemma}

\begin{proof}
Use the resolvent identity $R_A(z)=R_B(z)\bigl[I-(A-B)R_B(z)\bigr]^{-1}$ and Neumann
series when $\|(A-B)R_B(z)\|<1$, then integrate on~$\Gamma$ for the projection bounds.
\end{proof}

\subsection*{Applications to $L_{1}$ and $L_{3/2}$}

We apply Lemma~\ref{lem:resolvent-bridge} with
\[
(A,B)\in\bigl\{(L_{1},(L_{1})_{K}),\ (L_{3/2},(L_{3/2})_{k})\bigr\}.
\]
All norms below are operator norms on the indicated Hardy space.

If for some $z\in\mathbb{C}$ we have
\[
\varepsilon_{K}\,\bigl\|R_{(L_{1})_{K}}(z)\bigr\| \;<\; 1
\quad\text{or}\quad
\delta_{k}\,\bigl\|R_{(L_{3/2})_{k}}(z)\bigr\| \;<\; 1,
\]
then $z\in\rho(L_{1})$ or $z\in\rho(L_{3/2})$, respectively, and
\[
\|R_{L_{1}}(z)\|
\le \frac{\|R_{(L_{1})_{K}}(z)\|}{1-\varepsilon_K\|R_{(L_{1})_{K}}(z)\|},
\qquad
\|R_{L_{3/2}}(z)\|
\le \frac{\|R_{(L_{3/2})_{k}}(z)\|}{1-\delta_k\|R_{(L_{3/2})_{k}}(z)\|}.
\]
Thus validated upper bounds on the discretized resolvent certify **spectral gaps** for
$L_{1}$ or $L_{3/2}$.

Let $\Gamma$ be a simple closed contour containing a cluster of eigenvalues of the
discretization (but no other spectrum), and suppose
\[
\sup_{z\in\Gamma}\ \varepsilon_{K}\,\bigl\|R_{(L_{1})_{K}}(z)\bigr\|\ <\ 1
\quad\text{or}\quad
\sup_{z\in\Gamma}\ \delta_{k}\,\bigl\|R_{(L_{3/2})_{k}}(z)\bigr\|\ <\ 1.
\]
Then $\Gamma$ encloses exactly the same total algebraic multiplicity of eigenvalues for
$L_{1}$ (resp. $L_{3/2}$) as for the discretization. In particular, taking $\Gamma$ to be a
small circle around a simple eigenvalue of the discretization yields a validated enclosure
for the corresponding eigenvalue of $L_{1}$ or $L_{3/2}$.

Under the same hypothesis as in (B), the spectral projector onto the
eigenspace(s) inside $\Gamma$ satisfies
\[
\|P_{L_{1}}(\Gamma)-P_{(L_{1})_{K}}(\Gamma)\|
\ \le\ \frac{|\Gamma|}{2\pi}\ \sup_{z\in\Gamma}
\frac{\|R_{(L_{1})_{K}}(z)\|^2\,\varepsilon_{K}}{1-\varepsilon_{K}\|R_{(L_{1})_{K}}(z)\|},
\]
and the analogous bound holds with $L_{3/2}$ and $(L_{3/2})_{k}$, with $\delta_k$ in place
of $\varepsilon_K$. Hence eigenvectors and invariant subspaces (right/left) can be
approximated by integrating the discretized resolvent along~$\Gamma$ with an explicit
a posteriori error bound.

\medskip
\noindent\textbf{Remarks.}
(i) Since $(L_{1})_{K}$ and $(L_{3/2})_{k}$ have finite rank, their resolvent norms on
a circle $\Gamma$ can be computed (and rigorously enclosed) by any certified linear–algebra
routine (e.g., via Schur form and interval/ball arithmetic).  
(ii) Because the one–sided errors $\varepsilon_K,\delta_k$ decay like $(2/3)^{\cdot}$,
taking modest $K$ or $k$ often suffices to meet the small–gain condition
$\sup_{z\in\Gamma}\alpha(z)<1$.  
(iii) If $\Gamma$ surrounds the origin, note that the finite–rank maps may have $0$ in
their spectrum; choose $\Gamma$ avoiding~$0$ or use two–sided compressions if needed
(Thm.~\ref{thm:trunc-H2-correct}).

The resolvent identity and Riesz projections are standard; see, e.g., Kato, \emph{Perturbation Theory for Linear Operators}, Ch.~IV.


\section{Schur bridge on \texorpdfstring{$H^2(D_r)$}{H2(Dr)}}
\label{sec:schur-H2r}

Fix $r>0$ and work in the $N$–dimensional monomial subspace
$\mathcal V_N=\operatorname{span}\{(w-1)^n\}_{n=0}^{N-1}$.
In this basis the $H^2(D_r)$ Gram is diagonal,
\[
S_r=\operatorname{diag}(r^{2n})_{n=0}^{N-1},\qquad
C_r:=S_r^{1/2}=\operatorname{diag}(r^{n})_{n=0}^{N-1}.
\]
For a matrix $M\in\mathbb{C}^{N\times N}$ representing a linear map on $\mathcal V_N$,
the $H^2(D_r)$–operator norm is
\begin{equation}\label{eq:Hr-norm}
\|M\|_{(r)}:=\|\,C_r M C_r^{-1}\,\|_2 ,
\end{equation}
and the $S_r$–adjoint is $M^\sharp := S_r^{-1} M^\ast S_r$.

\section{Schur build in \texorpdfstring{$H^2(D_r)$}{H2(Dr)} and resolvent bounds}
\label{sec:schur-centered}

Fix \(r>0\) and the \(N\)-dimensional monomial subspace
\(\mathcal V_N=\mathrm{span}\{(w-1)^n\}_{n=0}^{N-1}\).
In this basis the \(H^2(D_r)\) Gram is diagonal
\[
S_r=\mathrm{diag}(r^{2n})_{n=0}^{N-1},\qquad C_r:=S_r^{1/2}=\mathrm{diag}(r^{n})_{n=0}^{N-1},
\]
and the induced \(H^2(D_r)\)-operator norm is
\[
\|M\|_{(r)}:=\|\,C_r\,M\,C_r^{-1}\,\|_2.
\]
For matrices \(X\), the \(S_r\)-adjoint is \(X^\sharp:=S_r^{-1}X^\ast S_r\).

\subsection*{Central procedure: build an \(H^2(D_r)\)-Schur pair}
Given a discretization \(A_k\in\mathbb C^{N\times N}\) of an endomorphism of \(H^2(D_r)\),
construct the Schur pair and its defects as follows.

\begin{enumerate}
\item \textbf{Whiten.} Set \(\widetilde A:=C_r\,A_k\,C_r^{-1}\).
Then \(\|\cdot\|_{(r)}\) becomes the Euclidean norm on \(\widetilde A\).

\item \textbf{Euclidean Schur.} Compute the (complex) Schur decomposition
\[
\widetilde A\,Q=Q\,T,
\]
with \(Q\) unitary and \(T\) upper (quasi-)triangular.

\item \textbf{Unwhiten the Schur vectors.} Define
\[
X:=C_r^{-1}Q,\qquad X^\sharp:=S_r^{-1}X^\ast S_r.
\]
Then \(A_k X = X T\) up to roundoff, and \(X^\sharp X \approx I\).

\item \textbf{Defects to be measured (cheaply).}
\[
\delta\ :=\ \|I-X^\sharp X\|_{(r)}=\|I-Q^\ast Q\|_2,\qquad
R\ :=\ A_kX-XT,\quad r_{\mathrm{res}}\ :=\ \|R\|_{(r)}=\|\widetilde A Q-QT\|_2.
\]
Set \(S_0:=XTX^\sharp\) and \(E:=A_k-S_0\).
\end{enumerate}

\medskip
All quantities above are obtained directly from \((Q,T)\) and diagonal scalings by \(C_r\).
In exact arithmetic, \(Q\) is unitary so \(\delta=0\); numerically \(\delta\) is at machine scale.

\begin{theorem}[From the Schur build to \(H^2(D_r)\) resolvent bounds]
\label{thm:schur-to-resolvent}
With the objects from the procedure above, let \(X=U H\) be the \(S_r\)-polar factorization
(\(U^\sharp U=I\), \(H=(X^\sharp X)^{1/2}\)). Then:
\begin{align*}
\quad&
\|X\|_{(r)}\le\sqrt{1+\delta},\\
\quad&
\|X^{-1}\|_{(r)}\le(1-\delta)^{-1/2},\\
\quad&
\|H\|_{(r)}\|H^{-1}\|_{(r)}\le\frac{1+\delta}{1-\delta}.\\
\quad&
 \ \|E\|_{(r)}\ \le\ r_{\mathrm{res}}\sqrt{1+\delta}\ +\ \|A_k\|_{(r)}\,\delta\ .\\
\quad&
\|(zI-S_0)^{-1}\|_{(r)}\ \le\ \frac{1+\delta}{1-\delta}\,\|(zI-T)^{-1}\|_2,\\
\end{align*}
And if 
\[
\alpha_r(z):=\|(zI-S_0)^{-1}\|_{(r)}\,\|E\|_{(r)}<1\
\] 
we have
\[
 \|(zI-A_k)^{-1}\|_{(r)}\ \le\
\frac{\dfrac{1+\delta}{1-\delta}\,\|(zI-T)^{-1}\|_2}{\,1-\alpha_r(z)\,}\ 
\]
\end{theorem}

\begin{proof}[Expanded proof]
Recall $\|M\|_{(r)}:=\|C_r M C_r^{-1}\|_2$ with $C_r=\textrm{diag}(r^n)$ and $X^\sharp=S_r^{-1}X^\ast S_r$.

\emph{(1) Orthogonality defect.}
If $\|I-X^\sharp X\|_{(r)}=\delta<1$, then (Loewner order)
\[
(1-\delta)I\preceq X^\sharp X\preceq(1+\delta)I,
\]
so
\[
\|X\|_{(r)}^2=\|X^\sharp X\|_{(r)}\le 1+\delta,\qquad
\|X^{-1}\|_{(r)}^{-2}=\|(X^\sharp X)^{-1}\|_{(r)}^{-1}\ge 1-\delta,
\]
hence
\[
\|X\|_{(r)}\le\sqrt{1+\delta},\qquad \|X^{-1}\|_{(r)}\le(1-\delta)^{-1/2}.
\]
With the $S_r$–polar $X=UH$ ($U^\sharp U=I$, $H=(X^\sharp X)^{1/2}$) we get
\[
\|H\|_{(r)}\|H^{-1}\|_{(r)}\le \frac{1+\delta}{1-\delta}.
\]

\emph{(2) Similarity defect (add \& subtract $A_kXX^\sharp$).}
Define $S_0:=XTX^\sharp$ and $E:=A_k-S_0$. Then
\[
E= \underbrace{A_kXX^\sharp - XTX^\sharp}_{(A_kX-XT)X^\sharp}\ +\ \underbrace{A_k - A_kXX^\sharp}_{A_k(I-XX^\sharp)}.
\]
Taking $\|\cdot\|_{(r)}$ gives
\[
\|E\|_{(r)}\le \|A_kX-XT\|_{(r)}\,\|X^\sharp\|_{(r)} + \|A_k\|_{(r)}\,\|I-XX^\sharp\|_{(r)}.
\]
Since $\|X^\sharp\|_{(r)}=\|X\|_{(r)}\le\sqrt{1+\delta}$ and (see below) $\|I-XX^\sharp\|_{(r)}=\|I-X^\sharp X\|_{(r)}=\delta$, we obtain
\[
\boxed{\quad \|E\|_{(r)}\ \le\ r_{\mathrm{res}}\sqrt{1+\delta}\ +\ \|A_k\|_{(r)}\,\delta,\quad}
\]
with $r_{\mathrm{res}}:=\|A_kX-XT\|_{(r)}$.

\emph{Why $\|I-XX^\sharp\|_{(r)}=\|I-X^\sharp X\|_{(r)}$.}
In whitened coordinates $\widetilde X:=C_r X C_r^{-1}$,
\[
\|I-XX^\sharp\|_{(r)}=\|I-\widetilde X\widetilde X^\ast\|_2,\qquad
\|I-X^\sharp X\|_{(r)}=\|I-\widetilde X^\ast\widetilde X\|_2.
\]
The nonzero eigenvalues of $\widetilde X\widetilde X^\ast$ and $\widetilde X^\ast\widetilde X$ are the same ($=\sigma_j(\widetilde X)^2$), so the eigenvalues of $I-$those are both $\{1-\sigma_j^2\}$; hence the spectral norms coincide.

\emph{(3) Resolvent comparison.}
With $X=UH$ we have $S_0=U(HTH)U^\sharp$, and $S_r$–unitary conjugation preserves $\|\cdot\|_{(r)}$, so
\[
\|(zI-S_0)^{-1}\|_{(r)}=\|(zI-HTH)^{-1}\|_{(r)}\le \|H\|_{(r)}\|H^{-1}\|_{(r)}\,\|(zI-T)^{-1}\|_2
\le \frac{1+\delta}{1-\delta}\,\|(zI-T)^{-1}\|_2.
\]

\emph{(4) Resolvent of $A_k$.}
Write
\[
zI-A_k=\big(I-(zI-S_0)^{-1}E\big)\,(zI-S_0).
\]
If $\alpha_r(z):=\|(zI-S_0)^{-1}\|_{(r)}\,\|E\|_{(r)}<1$, then Neumann’s lemma gives
\[
\|(zI-A_k)^{-1}\|_{(r)}\le \frac{\|(zI-S_0)^{-1}\|_{(r)}}{1-\alpha_r(z)}
\le \frac{\dfrac{1+\delta}{1-\delta}\,\|(zI-T)^{-1}\|_2}{1-\alpha_r(z)}.
\]
\end{proof}

\section{Newton--Kantorovich framework on a single Hardy space and biorthogonality}
\label{sec:NK-single}

Fix \(r\in\{1,\tfrac32\}\) and let \(L_r\) denote one of the endomorphisms
\[
L_1:=S\,L: H^2(D_1)\to H^2(D_1),
\qquad
L_{3/2}:=L\,R: H^2(D_{3/2})\to H^2(D_{3/2}),
\]
so that \(\|L_r\|\le C_2\) and any Galerkin compression \(A_k\) satisfies
\(\|A_k\|_{(r)}\le \|L_r\|\le C_2\) (Remark~\ref{rem:Ak-norms} below).
We assume \(\lambda\in\sigma(L_r)\setminus\{0\}\) is \emph{simple}.

\subsection*{Nonlinear map and biorthogonal normalization}
For a bounded \(M:H^2(D_r)\to H^2(D_r)\) define
\[
F_M(z,v)\ :=\ \big(Mv-zv,\ \ell(v)-1\big),\qquad (z,v)\in\mathbb C\times H^2(D_r),
\]
with a fixed linear functional \(\ell\in (H^2(D_r))^\ast\) used to pin phase/scale.
We take \(\ell\) from the \emph{discretization}: if \(A_k v_A=\lambda_A v_A\) and
\(u_A^\sharp A_k=\lambda_A u_A^\sharp\) are right/left eigenvectors, normalize by
\(\langle u_A,v_A\rangle_r=1\) and set \(\ell(\cdot)=\langle u_A,\cdot\rangle_r\).
Thus \(F_{A_k}(\lambda_A,v_A)=0\) and \(\ell(v_A)=1\).

\subsection*{Exact inverse of the linearization at a simple eigenpair}
Let \(P_A:=v_A\otimes u_A\) and \(Q_A:=I-P_A\).
Then for the derivative at \((\lambda_A,v_A)\),
\[
DF_{A_k}(\lambda_A,v_A)[\delta z,\delta v]
=\big((A_k-\lambda_A)\delta v-\delta z\,v_A,\ \langle u_A,\delta v\rangle_r\big).
\]
If \(\lambda_A\) is simple, the restriction
\((A_k-\lambda_A)|_{Q_A}:Q_AH^2(D_r)\to Q_AH^2(D_r)\) is invertible and
\[
\boxed{\quad
DF_{A_k}(\lambda_A,v_A)^{-1}(a,b)
= \big(\ -\langle u_A,a\rangle_r,\ \ b\,v_A + \underbrace{(A_k-\lambda_A)|_{Q_A}^{-1}Q_A a}_{=:~S_A a}\ \big).
\quad}
\]
Hence, with the product norm \(\|(a,b)\|:=\|a\|_{(r)}+|b|\),
\[
\|DF_{A_k}^{-1}\|\ \le\ \max\big\{\ \|u_A\|_{(r)}+\|S_A\|\ ,\ \|v_A\|_{(r)}\ \big\}.
\]
If \(\Gamma\) is a circle centered at \(\lambda_A\) that contains no other spectrum of \(A_k\),
then the \emph{reduced resolvent} obeys
\[
S_A=\frac{1}{2\pi i}\int_\Gamma \frac{1}{z-\lambda_A}\,(zI-A_k)^{-1}\,dz,
\qquad
\Rightarrow\quad
\boxed{\ \ \|S_A\|\ \le\ \sup_{z\in\Gamma}\|(zI-A_k)^{-1}\|_{(r)}\ .\ }
\]

\subsection*{Newton--Kantorovich transfer to $L_r$ (same space)}
Let \((\lambda,v)\) be the (normalized) simple eigenpair of \(L_r\) with
\(\ell(v)=1\) and \(L_r v=\lambda v\). Consider the Newton maps
\(N_{A_k}:=(\mathrm{id}-DF_{A_k}^{-1}F_{A_k})\) and \(N_{L_r}\).
If \(N_{A_k}\) is a contraction on a ball around \((\lambda_A,v_A)\) with constant \(\kappa<1\),
then (telescoping at the fixed points and using \(F_{L_r}(\lambda,v)=0\)) we get
\[
\big\|(\lambda_A,v_A)-(\lambda,v)\big\|
\ \le\ \frac{1}{1-\kappa}\ \|DF_{A_k}^{-1}\|\ \|(A_k-L_r)v\|_{(r)}.
\]
Thus, using \(\|(A_k-L_r)v\|\le \|A_k-L_r\|\,\|v\|\),
\[
\boxed{\quad
\big\|(\lambda_A,v_A)-(\lambda,v)\big\|
\ \le\ \frac{\|DF_{A_k}^{-1}\|}{1-\kappa}\ \|A_k-L_r\|_{(r)}\ \|v\|_{(r)}.
\quad}
\]

Now, 

In particular, if \(A_k\) is a Galerkin compression produced by range/domain truncation,
then (Corollary~\ref{cor:one-sided-approxs})
\[
\|A_k-L_1\|_{(1)}\ \le\ C_2\Bigl(\tfrac{2}{3}\Bigr)^{K+1},
\qquad
\|A_k-L_{3/2}\|_{(3/2)}\ \le\ C_2\Bigl(\tfrac{2}{3}\Bigr)^{k+1},
\]
and the right-hand side decays geometrically with \(K\) (or \(k\)).

\begin{remark}[Using resolvent bounds for \(\|DF_{A_k}^{-1}\|\)]
If \(\Gamma\) is a circle around \(\lambda_A\) with no other eigenvalues of \(A_k\), then
\[
\|DF_{A_k}^{-1}\|
\ \le\ \max\Big\{\ \|u_A\|_{(r)} + \sup_{z\in\Gamma}\|(zI-A_k)^{-1}\|_{(r)}\ ,\ \|v_A\|_{(r)}\ \Big\}.
\]
Both \(\|v_A\|\) and \(\|u_A\|\) are computable from the Schur build in \(H^2(D_r)\),
and the resolvent supremum on \(\Gamma\) admits certified bounds via
Theorem~\ref{thm:schur-to-resolvent}.
\end{remark}

\begin{remark}[Norm control of the discretizations]\label{rem:Ak-norms}
If \(A_k\) is the \(H^2(D_r)\)-orthogonal compression of \(L_r\) onto
\(\mathcal V_N=\mathrm{span}\{(w-1)^n\}_{n=0}^{N-1}\), then
\[
\|A_k\|_{(r)}\ \le\ \|L_r\|\ \le\ C_2,\qquad r\in\{1,\tfrac32\}.
\]
Indeed, \(\|L_1\|\le \|S\|\,\|L\|=C_2\) and \(\|L_{3/2}\|\le \|L\|\,\|R\|=C_2\)
with \(\|R\|=\|S\|=1\), and orthogonal compression does not increase the operator norm.
\end{remark}

\paragraph{Replacing $\|v\|_{(r)}$ in the NK bound.}
In the single–space framework of Section~\ref{sec:NK-single} we derived
\[
\big\|(\lambda_A,v_A)-(\lambda,v)\big\|
\ \le\ \frac{\|DF_{A_k}^{-1}\|}{1-\kappa}\ \|A_k-L_r\|_{(r)}\ \|v\|_{(r)}.
\tag{$\ast$}
\]
There are two convenient ways to control $\|v\|_{(r)}$:

\smallskip
\noindent\emph{(A) Same–space normalization (recommended).}
Normalize the true eigenvector in the working space: $\|v\|_{(r)}=1$.
Then $(\ast)$ becomes
\[
\big\|(\lambda_A,v_A)-(\lambda,v)\big\|
\ \le\ \frac{\|DF_{A_k}^{-1}\|}{1-\kappa}\ \|A_k-L_r\|_{(r)}.
\]

\smallskip
\noindent\emph{(B) Weak-to-strong conversion (when the target normalization is in $H^2(D_1)$).}
Suppose $r=\tfrac32$ but you keep the weak normalization $\|v\|_{H^2(D_1)}=1$.
By Proposition~\ref{prop:strong-from-weak},
\[
\|v\|_{(3/2)}=\|v\|_{H^2(D_{3/2})}\ \le\ \frac{C_2}{|\lambda|}\,\|v\|_{H^2(D_1)}
\ =\ \frac{C_2}{|\lambda|}.
\]
Plugging this into $(\ast)$ with $r=\tfrac32$ yields the one–norm estimate
\[
\boxed{\qquad
\big\|(\lambda_A,v_A)-(\lambda,v)\big\|
\ \le\ \frac{C_2}{(1-\kappa)\,|\lambda|}\ \|DF_{A_k}^{-1}\|\ \|A_k-L_{3/2}\|_{(3/2)}.
\qquad}
\]
(The same argument shows $\|v\|_{(3/2)}\le (C_2/|\lambda|)\|v\|_{(1)}$ for any scaling.)



\section{Polynomials of operators on \texorpdfstring{$H^2(D_r)$}{H2(Dr)}}
\label{sec:poly-single}

Fix \(r\in\{1,\tfrac32\}\) and write \(L_r\in\mathcal B(H^2(D_r))\) for either endomorphism
\(L_1=SL\) or \(L_{3/2}=LR\) (so \(\|L_r\|\le C_2\)).
Let \(A_k\) be a Galerkin compression on \(\mathcal V_N=\mathrm{span}\{(w-1)^n\}_{n=0}^{N-1}\),
and set the single–space discretization error
\[
\varepsilon_r\ :=\ \|L_r-A_k\|_{(r)}.
\]

\subsection{Polynomial perturbation in one norm}
Let \(p(z)=\sum_{j=1}^{d} a_j z^j\) with \(p(0)=0\).
A telescoping expansion gives
\[
p(L_r)-p(A_k)
=\sum_{j=1}^{d} a_j \sum_{\ell=0}^{j-1}
A_k^{\,\ell}\,(L_r-A_k)\,L_r^{\,j-1-\ell},
\]
hence
\begin{equation}\label{eq:poly-perturb-single}
\boxed{\quad
\|p(L_r)-p(A_k)\|_{(r)}
\ \le\ \varepsilon_r\,
\underbrace{\sum_{j=1}^{d} |a_j|\sum_{\ell=0}^{j-1}
\|A_k\|_{(r)}^{\ell}\,\|L_r\|_{(r)}^{\,j-1-\ell}}_{=:~\mathfrak C_r(p;A_k,L_r)}.
\quad}
\end{equation}
Crude but convenient: using \(\|A_k\|_{(r)},\|L_r\|_{(r)}\le C_2\),
\begin{equation}\label{eq:poly-perturb-C2}
\|p(L_r)-p(A_k)\|_{(r)}
\ \le\ \varepsilon_r \sum_{j=1}^{d} |a_j|\,j\,C_2^{\,j-1}.
\end{equation}

\subsection{Resolvent transfer for polynomials (same space)}
\begin{lemma}[Single–space resolvent transfer]\label{lem:resolvent-transfer-single}
Fix a contour \(\Gamma\subset\mathbb{C}\). Let
\[
M_r\ :=\ \sup_{z\in\Gamma}\|(zI-p(A_k))^{-1}\|_{(r)},
\qquad
\varepsilon_{p,r}\ :=\ \|p(L_r)-p(A_k)\|_{(r)}.
\]
If \(\varepsilon_{p,r}\,M_r<1\), then \(\Gamma\subset\rho(p(L_r))\) and
\[
\sup_{z\in\Gamma}\|(zI-p(L_r))^{-1}\|_{(r)}
\ \le\ \frac{M_r}{\,1-\varepsilon_{p,r} M_r\,}.
\]
\end{lemma}

\begin{proof}[Sketch]
Resolvent identity:
\(z-p(L_r)=\big[I-(p(L_r)-p(A_k))(z-p(A_k))^{-1}\big]\,(z-p(A_k))\).
Apply Neumann’s lemma when \(\varepsilon_{p,r}M_r<1\).
\end{proof}

\subsection{Computing \texorpdfstring{$M_r$}{Mr} from a Schur build of \texorpdfstring{$A_k$}{Ak}}
Use the \(H^2(D_r)\) Schur procedure of Section~\ref{sec:schur-centered}:
whiten \(A_k\) by \(C_r\), compute the Schur pair \(\widetilde A Q=QT\), set \(X=C_r^{-1}Q\),
\(\delta=\|I-X^\sharp X\|_{(r)}\) and \(r_{\mathrm{res}}=\|A_kX-XT\|_{(r)}\).
Since \(p\) is a polynomial,
\[
p(S_0)=X\,p(T)\,X^\sharp,\qquad S_0:=XTX^\sharp.
\]
Define \(E_p:=p(A_k)-p(S_0)\). Arguing as in Theorem~\ref{thm:schur-to-resolvent},
\[
\|(zI-p(S_0))^{-1}\|_{(r)}\ \le\ \frac{1+\delta}{1-\delta}\ \|(zI-p(T))^{-1}\|_2,
\]
and, by the same add–subtract trick used for \(E\),
\begin{equation}\label{eq:poly-sim-defect}
\|E_p\|_{(r)}\ \le\ \sqrt{1+\delta}\,\sum_{j=1}^d |a_j|
\sum_{\ell=0}^{j-1} \|A_k\|_{(r)}^{\ell}\, r_{\mathrm{res}}\ \|T\|_2^{\,j-1-\ell}
\ +\ \delta\,\|p(A_k)\|_{(r)}.
\end{equation}
If \(\alpha_p(z):=\|(zI-p(S_0))^{-1}\|_{(r)}\,\|E_p\|_{(r)}<1\), then
\[
\boxed{\quad
\|(zI-p(A_k))^{-1}\|_{(r)}
\ \le\ \frac{\dfrac{1+\delta}{1-\delta}\,\|(zI-p(T))^{-1}\|_2}{1-\alpha_p(z)}\ =:\ M_r(z),
\quad}
\]
and \(M_r:=\sup_{z\in\Gamma}M_r(z)\) is a certified bound for Lemma~\ref{lem:resolvent-transfer-single}.
In exact arithmetic \(Q\) is unitary in the whitened metric (\(\delta=0\), \(r_{\mathrm{res}}=0\)) and
\(\|(zI-p(A_k))^{-1}\|_{(r)}=\|(zI-p(T))^{-1}\|_2\).

\subsection{Deflation polynomials and isolation}
Let \(\{\lambda_i\}\) be the Schur–ordered eigenvalues of \(A_k\) (in \(H^2(D_r)\)).
Suppose \(\{\widehat\lambda_i\}_{i<j}\) are certified. Define the deflation polynomial
\[
P_j(z):=\prod_{i=1}^{j-1}\Bigl(1-\frac{z}{\widehat\lambda_i}\Bigr),\qquad
\alpha_j:=\frac{1}{P_j(\lambda_j)},\qquad
p(z):=\bigl(\alpha_j P_j(z)\bigr)^q,
\]
with the minimal integer \(q\ge1\) that yields the desired order of vanishing at the
already certified set while keeping \(p(\lambda_j)\approx1\).
Then the spectrum of \(p(A_k)\) has one point near \(1\) (image of \(\lambda_j\))
and the rest contracted toward \(0\). Any small circle \(\Gamma=\{z:|z-1|=r\}\) that
excludes the other images isolates the target.

\subsection{Projectors commute with polynomials}
If \(\Gamma\) encloses a single eigenvalue of \(A_k\) (or \(L_r\)), then the Riesz
projector is invariant under polynomial change:
\[
\Pi=\frac{1}{2\pi i}\int_\Gamma (z-A)^{-1}\,dz
=\frac{1}{2\pi i}\int_{p(\Gamma)} (w-p(A))^{-1}\,dw,
\]
provided \(p\) is injective on the enclosed component. Thus one may compute
\(\widehat\Pi_j\) from \(p(A_k)\) and transfer it to \(p(L_r)\) via
Lemma~\ref{lem:resolvent-transfer-single}; mapping back through the local inverse of
\(p\) yields the certified eigenvalue and projector for \(L_r\).

\subsection{Practical pipeline on \texorpdfstring{$H^2(D_r)$}{H2(Dr)}}
\begin{enumerate}
\item \textbf{Schur build for \(A_k\).} Whiten by \(C_r\), compute \(\widetilde A Q=QT\),
set \(X=C_r^{-1}Q\), measure \(\delta, r_{\mathrm{res}}\).
\item \textbf{Choose \(p\).} Deflate previously certified modes; take minimal \(q\).
\item \textbf{Bound \(M_r\).} Use \eqref{eq:poly-sim-defect} and the resolvent bound with \(p(T)\)
to certify \(M_r=\sup_{z\in\Gamma}\|(zI-p(A_k))^{-1}\|_{(r)}\).
\item \textbf{Bound \(\varepsilon_{p,r}\).} Use \eqref{eq:poly-perturb-single} (or the crude
\eqref{eq:poly-perturb-C2} if desired) with \(\varepsilon_r=\|L_r-A_k\|_{(r)}\).
\item \textbf{Transfer.} Check \(\varepsilon_{p,r} M_r<1\). Then
\[
\sup_{z\in\Gamma}\|(zI-p(L_r))^{-1}\|_{(r)}\ \le\ \frac{M_r}{1-\varepsilon_{p,r}M_r},
\]
so \(\Gamma\subset\rho(p(L_r))\) and the projector is well defined with matching multiplicity.
\item \textbf{Back–map.} Use the local inverse of \(p\) at the image of \(\lambda_j\) to enclose
\(\lambda_j\in\sigma(L_r)\) and recover the projector on \(H^2(D_r)\).
\end{enumerate}

\begin{remark}[Constants you actually need]
All inputs are single–space: \(C_2\) (for crude bounds), the discretization error
\(\varepsilon_r=\|L_r-A_k\|_{(r)}\) (from your one–sided truncation bounds),
the Schur defects \(\delta, r_{\mathrm{res}}\), and the triangular resolvent
\(\|(zI-p(T))^{-1}\|_2\). No cross–norms or embeddings are required.
\end{remark}




\section{Outlook}
This framework couples Hardy-space analysis and validated numerics.
The same operator \(L_s\) for complex \(s\) yields
\(\det(I-L_s)^{-1}=\zeta(2s)\),
forming the analytic basis for a companion paper.

\begin{remark}[Connection with the Markov and Lagrange spectra]
The parameter \(s\) in the Gauss--Kuzmin--Wirsing operator
\[
(L_s f)(w)=\sum_{n\ge1}\frac{1}{(w+n)^{2s}}f\!\left(\frac{1}{w+n}\right)
\]
acts as a thermodynamic deformation of the Gauss map potential.
At \(s=1\), the operator describes the statistical law of continued fractions
and governs the convergence of the empirical distribution functions \(F_n\).
As \(s\) decreases towards \(1/2\), the leading eigenvalue
\(\lambda_1(s)=e^{P(2s)}\)
traces the topological pressure \(P(2s)\),
whose zero determines the Hausdorff dimension
of sets of numbers with bounded continued fractions:
\[
P(2D)=0,\qquad D\approx0.53128.
\]
This critical value \(D\) marks the onset of the
Markov--Lagrange spectrum.
Hence the same analytic family \(L_s\)
interpolates between the probabilistic regime (\(s=1\))
and the geometric--Diophantine regime (\(s=D\)),
linking the Gauss problem with the classical spectra of Diophantine approximation.
\end{remark}

\section{Appendix: Pseudospectral approach to the enclosure of eigenvalues strong-weak framework}

Recall that $L:H^2(D_1)\to H^2(D_{3/2})$ is bounded with
\[
\|L\|_{H^2(D_1)\to H^2(D_{3/2})}\ \le\ C_2.
\]
(For instance, one can take $C_2$ from Theorem~\ref{thm:L1-H2-bound}.)

\begin{proposition}[Strong norm from weak norm for eigenvectors]
\label{prop:strong-from-weak}
Let $(\lambda,\varphi)$ be a (nonzero) eigenpair of $L$ with
$\varphi\in H^2(D_1)\cap H^2(D_{3/2})$ and $L\varphi=\lambda\varphi$.
Then
\[
\boxed{\qquad
\|\varphi\|_{H^2(D_{3/2})}
\ \le\ \frac{C_2}{|\lambda|}\,\|\varphi\|_{H^2(D_1)}.
\qquad}
\]
In particular, under the normalization $\|\varphi\|_{H^2(D_1)}=1$,
\[
\|\varphi\|_{H^2(D_{3/2})}\ \le\ \frac{C_2}{|\lambda|}.
\]
\end{proposition}

\begin{proof}
Since $L\varphi=\lambda\varphi$ and $L:H^2(D_1)\to H^2(D_{3/2})$,
\[
\|\varphi\|_{H^2(D_{3/2})}
=\frac{1}{|\lambda|}\,\|L\varphi\|_{H^2(D_{3/2})}
\le \frac{C_2}{|\lambda|}\,\|\varphi\|_{H^2(D_1)}.
\]
\end{proof}

We now combine Proposition~\ref{prop:strong-from-weak} with a resolvent
estimate to obtain an exclosure for $\sigma(L)$ from a finite-rank
approximation $L_k$.

\begin{proposition}\label{prop:pseudo-exclosure-correct}
Let $L_k$ be a finite-rank operator such that
\[
\|L-L_k\|_{H^{2}(D_{3/2})\to H^{2}(D_{1})} \le \eta.
\]
If $\lambda\in\sigma(L)$ and $\lambda\notin\sigma(L_k)$, then
\[
\bigl\|(\lambda - L_k)^{-1}\bigr\|_{H^{2}(D_{1})\to H^{2}(D_{1})}
\ \ge\ \frac{|\lambda|}{C_2\,\eta}.
\]
Equivalently,
\[
\sigma(L)\ \subset\
\left\{\, z\in\mathbb C\ \middle|\ 
\bigl\|(z - L_k)^{-1}\bigr\|_{H^{2}(D_{1})\to H^{2}(D_{1})}
\ \ge\ \frac{|z|}{C_2\,\eta}\right\}.
\]
Consequently, any region where
$\|(z-L_k)^{-1}\|_{H^{2}(D_{1})}\le 1/\delta$ with $\delta>\frac{C_2\,\eta}{|z|}$
is free of eigenvalues of $L$.
\end{proposition}

\begin{proof}
Let $(\lambda,\varphi)$ be an eigenpair with $\varphi\neq 0$ and
$\varphi\in H^2(D_1)\cap H^2(D_{3/2})$. Then
\[
(\lambda - L_k)\varphi = (L - L_k)\varphi,
\]
so, assuming $\lambda\notin\sigma(L_k)$,
\[
\varphi \;=\; (\lambda - L_k)^{-1}(L - L_k)\varphi.
\]
Taking the $H^2(D_1)$-norm and using
$\|L-L_k\|_{H^{2}(D_{3/2})\to H^{2}(D_{1})}\le\eta$ gives
\[
\|\varphi\|_{H^2(D_1)}
\;\le\;
\bigl\|(\lambda - L_k)^{-1}\bigr\|_{H^{2}(D_{1})\to H^{2}(D_{1})}\;
\eta\;\|\varphi\|_{H^2(D_{3/2})}.
\]
Apply Proposition~\ref{prop:strong-from-weak} to bound
$\|\varphi\|_{H^2(D_{3/2})}\le (C_2/|\lambda|)\,\|\varphi\|_{H^2(D_1)}$ and cancel
$\|\varphi\|_{H^2(D_1)}\neq 0$ to obtain
\(
1 \le \bigl\|(\lambda - L_k)^{-1}\bigr\|\;\eta\; (C_2/|\lambda|),
\)
i.e.
\(
\bigl\|(\lambda - L_k)^{-1}\bigr\|
\ge |\lambda|/(C_2\,\eta).
\)
The set inclusion and the exclosure statement follow by contraposition.
\end{proof}

\begin{lemma}[Energy inequality for the resolvent equation]
\label{lem:energy-ineq}
Let $(\mathcal B_s,\|\cdot\|_s)\hookrightarrow(\mathcal B_w,\|\cdot\|_w)$ be a Banach couple
and suppose $P:\mathcal B_s\to\mathcal B_s$ satisfies the strong Lasota--Yorke estimate
\[
\|Pu\|_s \le a\,\|u\|_s + b\,\|u\|_w, \qquad a,b\ge0.
\]
Fix $z\in\rho(P)$ and let $u\in\mathcal B_s$ solve the resolvent equation
\[
(zI-P)u = f, \qquad f\in\mathcal B_s.
\]
Then the solution obeys the \emph{energy inequality}
\[
\boxed{\ (|z|-a)\,\|u\|_s \;\le\; b\,\|u\|_w \;+\; \|f\|_s.\ }
\]
In particular, the strong norm of $u$ is controlled by its weak norm and the data $f$.
\end{lemma}

\begin{proof}
From $(zI-P)u=f$ we have
\[
|z|\,\|u\|_s \le \|Pu\|_s + \|f\|_s.
\]
Applying the Lasota--Yorke bound gives
\[
|z|\,\|u\|_s \le a\,\|u\|_s + b\,\|u\|_w + \|f\|_s,
\]
and rearranging yields the claimed inequality.
\end{proof}

\begin{proposition}
\label{prop:weak-to-strong-bridge}
Let $(\mathcal B_s,\|\cdot\|_s)\hookrightarrow(\mathcal B_w,\|\cdot\|_w)$ with continuous embedding
$i:\mathcal B_s\to\mathcal B_w$ of norm $E_{s\to w}=\|i\|_{s\to w}$.
Assume $P,P_k:\mathcal B_s\to\mathcal B_s$ satisfy the Lasota--Yorke bounds
\[
\|Pu\|_s \le a\,\|u\|_s + b\,\|u\|_w, \qquad
\|Pu\|_w \le M\,\|u\|_w,
\]
and let $\delta:=\|P-P_k\|_{s\to w}$.
Suppose for $z\notin\sigma(P_k)$ we have a weak resolvent bound
\[
M_k(z):=\|(zI-P_k)^{-1}\|_{w\to w} < \infty.
\]
If
\[
|z| > a + b\,M_k(z)\,\delta,
\]
then $z\notin\sigma(P)$ and the resolvent $R(z)=(zI-P)^{-1}$
satisfies the quantitative estimate
\[
\boxed{\;
\|R(z)\|_{s\to s}
\ \le\
\frac{\,b\,M_k(z)\,E_{s\to w} + 1\,}{\,|z|-a - b\,M_k(z)\,\delta\,}.
\;}
\]
\end{proposition}

\begin{proof}
Let $R_k(z)=(zI-P_k)^{-1}$ and fix $f\in\mathcal B_s$ with $\|f\|_s=1$.
Set $u=R(z)f$, so that $(zI-P)u=f$ and hence
\[
u = R_k(z)f + R_k(z)(P-P_k)u.
\]
Taking weak norms gives
\[
\|u\|_w
 \le M_k(z)\,\|f\|_w + M_k(z)\,\|(P-P_k)u\|_w
 \le M_k(z)\,E_{s\to w} + M_k(z)\,\delta\,\|u\|_s.
\]
Apply Lemma~\ref{lem:energy-ineq} to $(zI-P)u=f$:
\[
(|z|-a)\,\|u\|_s \le b\,\|u\|_w + 1.
\]
Combining the two inequalities yields
\[
(|z|-a)\,\|u\|_s \le b M_k(z)\,E_{s\to w} + b M_k(z)\,\delta\,\|u\|_s + 1,
\]
collecting terms,
\[
\big(|z|-a - b\,M_k(z)\,\delta\big)\,\|u\|_s
   \le b\,M_k(z)\,E_{s\to w} + 1.
\]
If $|z|>a+b\,M_k(z)\,\delta$, divide through to obtain the bound;
taking the supremum over $\|f\|_s=1$ gives the stated estimate.
\end{proof}




\bibliographystyle{plain}
\bibliography{refs}
\end{document}
